% =========================================================================
% Chapter 4: Results and Discussion
%
% Student Guide:
% - Shows your results (tables, graphs) and explains what they mean.
% - Link your findings back to your objectives and to earlier work.
% =========================================================================

\chapter{Results and Discussion}
\label{ch:results}

The Results and Discussion chapter is the heart of your thesis, where raw data is transformed into meaningful insights. It presents your findings, explains their significance, and connects them to your research aims, hypotheses, and existing literature. Unlike the methodology or literature review, this chapter focuses on evidence and interpretation, demonstrating the value of your research in Computer Science and Engineering.

\section{Presenting Results}
\label{sec:presenting-results}

Results should be presented in a clear, logical, and structured manner. Use tables, figures, and graphs to illustrate key outcomes, ensuring they are properly labeled and referenced. Report results objectively, without interpretation at this stage. Organize them in a way that best supports your study, such as:

\begin{enumerate}
  \item From most important to least important (e.g., from statistically most significant to least significant).
  \item By research questions or problems, presenting results that directly answer each question sequentially.
  \item By methods or experimental elements, grouping procedures with their corresponding results.
\end{enumerate}

Summarize data rather than overwhelming readers with raw outputs, and highlight findings that directly address your research objectives.

\section{Interpreting Results in Discussion}
\label{sec:interpreting-results}

After presenting the data, move into the discussion where you interpret and explain what the results mean. Relate them directly to your research aims and hypotheses, showing whether they confirm, extend, or challenge expectations. Compare your findings with prior studies and theoretical frameworks, highlighting similarities, differences, and possible reasons for discrepancies. This critical analysis situates your work within the broader academic context.

\section{Connecting to Research Objectives}
\label{sec:connecting-objectives}

Explicitly connect your discussion back to the objectives and research questions stated earlier in the thesis. For each objective, explain how your findings provide answers or insights. This alignment reinforces coherence and demonstrates that your study has achieved its intended goals.

\section{Addressing Limitations}
\label{sec:addressing-limitations}

A strong discussion acknowledges limitations. Reflect on constraints such as dataset size, simulation environment, or methodological assumptions that may have influenced the results. Transparency about limitations demonstrates academic integrity and provides direction for future research.

\section{Highlighting Contributions and Implications}
\label{sec:ch4-highlighting-contributions}

Discuss the significance of your findings. Explain how they advance theoretical understanding, improve practical systems, or have potential industrial or societal applications. For example, if your work improves algorithm efficiency, highlight its impact on real-world computing or communication systems.

\section{Structuring the Chapter}
\label{sec:structuring-results-chapter}

You may choose to:
\begin{itemize}
  \item Separate results and discussion into distinct sections, or
  \item Integrate them together, presenting each finding followed immediately by its interpretation.
\end{itemize}
The choice depends on your discipline and supervisor's guidance, but clarity and logical flow are essential.

The rest of this chapter provides some additional styles that you are going to use in your thesis report. Styles used in other parts of the document:

\section{Thesis Title Style}
\label{sec:thesis-title-style}

The ThesisTitle style is used on the first page and second page of the document to format the thesis title. In this \LaTeX{} template, the title is defined centrally in \texttt{thesis-config.tex} via the \texttt{\textbackslash thesisTitle} command. Room is kept for a title of up to four lines. Whether your title takes one line or four, nothing else on those two pages moves.

\section{Department School University Style}
\label{sec:dept-school-univ-style}

The Dept-College-Unv. style is applied on the first page and second page to format the departmental information, ensuring it appears centered and distinct from other text. In this \LaTeX{} template, institutional credentials are configured centrally in \texttt{thesis-config.tex}.
